%  American Terawatt — SMP Level 1, single-line diagram
%
%  Structure:
%    amwatt_20261001.tex         <- this file
%    ../../template/theme        <- shared beamer theme (atstyle.sty, fonts, img)
%    .latexmkrc                  <- selects xelatex
%    fig/*.pdf                   <- the figures, as finished PDF drawings
%
%  Compile in this folder:
%    latexmk -xelatex amwatt_20261001.tex
%
%  The figures are included as PDFs with \includegraphics -- no TikZ
%  drawing and no Glyphwright symbol library in this document. Anything
%  laid over a figure (power-flow arrows, section bands, voltage
%  annotations) is positioned in FRACTIONS OF THE INCLUDED IMAGE: the
%  image is one TikZ node named `fig', and (f) maps its bounding box to
%  the unit square, so (f) of (0,0) is the lower left corner and (f) of
%  (1,1) the upper right. Replacing a PDF with a differently cropped one
%  therefore only needs those fractions adjusted, nothing else.
% =============================================================
\documentclass[aspectratio=169,11pt]{beamer}
% The beamer theme lives in the shared template folder; \atthemepath
% tells it where its own fonts/ and img/ are, relative to this file.
\providecommand{\atthemepath}{../../template/theme/}
\usepackage{../../template/theme/atstyle}

\usetikzlibrary{arrows.meta,calc}

\title{HVDC Design}
\subtitle{AmWatt Team Presentation}
\author{American Terawatt}
\date{\today}

% ---- figure helpers -----------------------------------------
% \figbox{<width>}{<name>} puts fig/<name>.pdf into a TikZ node called
% `fig' and names its four bounding-box corners f-sw, f-se, f-nw, f-ne,
% so that anything laid over the figure can be placed in fractions of
% its width and height, e.g. ($(f-sw)!0.38!(f-se)$).
\newcommand{\figfolder}{fig/}
\newcommand{\figbox}[2]{%
  \node[inner sep=0] (fig) {\includegraphics[width=#1]{\figfolder#2}};%
  \coordinate (f-sw) at (fig.south west);%
  \coordinate (f-se) at (fig.south east);%
  \coordinate (f-nw) at (fig.north west);%
  \coordinate (f-ne) at (fig.north east);%
}

\definecolor{gwvoltage}{RGB}{31,95,191}

\tikzset{
  pflow/.style={-{Latex[length=2mm]}, violet, line width=1.2pt},
  pboth/.style={pflow, {Latex[length=2mm]}-{Latex[length=2mm]}},
  pac2dc/.style={pflow, -{Latex[length=2mm]}},
  pdc2ac/.style={pflow, {Latex[length=2mm]}-},
  volt/.style={-{Latex[length=2mm]}, line width=0.8pt, draw=gwvoltage},
  gnd/.style={draw=gwvoltage, line width=0.6pt},
  band/.style={fill opacity=0.15, rounded corners=1pt, draw=none},
  modelbl/.style={font=\small\bfseries},
}

  % \vscmode{node name}{position}{label}{arrow style}
  \def\vscmode#1#2#3#4{%
    \node[inner sep=0] (#1) at (#2)
      {\includegraphics[width=0.15\textwidth]{\figfolder vsc_conv}};
    \node[modelbl, anchor=south] at ($(#1.north) + (0,1.5mm)$) {#3};
    \draw[#4] ($(#1.south) + (-10mm,-3mm)$) -- node[below] {$P$}
              ($(#1.south) + (10mm,-3mm)$);
  }

\begin{document}

\titleslide

% ---- Single-line diagram ------------------------------------
\begin{frame}{System Overview}
  ORC generators feed a data centre through a point-to-point VSC HVDC link.
  \vspace{2.2em}

  \centering
  \begin{tikzpicture}
    \figbox{0.92\textwidth}{smp_level_1_gen_load}
    \draw[pflow] ([xshift=-6mm]fig.center) -- node[above] {$P$}
                 ([xshift=6mm]fig.center);
  \end{tikzpicture}
\end{frame}

% \band{colour}{from}{to}: a section band spanning two positions given as
% fractions of the figure width (defined outside the frame, since beamer
% does not accept macro parameters inside frame bodies).
\newcommand{\band}[3]{%
  \fill[band, #1] ([yshift=-1.5mm]$(f-sw)!#2!(f-se)$)
         rectangle ([yshift=1.5mm]$(f-nw)!#3!(f-ne)$);}

% ---- Single-line diagram with the supply sections -----------
\begin{frame}{System Overview}
  \centering
  The complete system consists of different elements, supplied by different
  companies.\\
  AmWatt is responsible for the integration of those elements.
  \vspace{1.6em}

  \centering
  \begin{tikzpicture}
    \figbox{0.92\textwidth}{smp_level_1_gen_load}
    % The boundaries sit midway between the elements of the drawing: left
    % busbar at 0.133, converter boxes at 0.286/0.373 and 0.647/0.734, right
    % busbar at 0.887.
    \band{atLime}{0.000}{0.123}   % generation
    \band{atGray}{0.125}{0.143}   % busbar
    \band{cyan}{0.145}{0.382}     % converter station
    \band{violet}{0.384}{0.636}   % d.c. line
    \band{cyan}{0.638}{0.877}     % converter station
    \band{atGray}{0.879}{0.897}   % busbar
    \band{orange}{0.899}{1.001}   % data centre
    \draw[pflow] ([xshift=-6mm]fig.center) -- node[above] {$P$} ([xshift=6mm]fig.center);
  \end{tikzpicture}

\end{frame}

% ---- VSC converter -------------------------------------------
\begin{frame}{VSC converter}
  \noindent The VSC converts the a.c.\ voltage $U_{ac}$ into the pole-to-pole
  d.c.\ voltage $U_{dc}$.
  \vspace{1.5em}

  \centering
  \begin{tikzpicture}
    \figbox{0.20\textwidth}{vsc_conv}
    % a.c. terminal: the lead leaves the left edge at mid height; its
    % voltage is drawn against earth below the figure
    \coordinate (acp) at ($(fig.west) + (0.6mm,0)$);
    \fill (acp) circle (0.6pt);
    \draw[volt] ($(acp)+(0,-1mm)$) -- ++(0,-12mm) coordinate (acg);
    \draw[gnd] ($(acg)+(-2.2mm,-0.6mm)$) -- ++(4.4mm,0);
    \draw[gnd] ($(acg)+(-1.3mm,-1.9mm)$) -- ++(2.6mm,0);
    \draw[gnd] ($(acg)+(-0.5mm,-3.2mm)$) -- ++(1.0mm,0);
    \node[anchor=east, color=gwvoltage] at ($(acp)!0.5!(acg) + (-1.5mm,0)$)
      {$U_{ac}$};
    % d.c. terminals: the two leads leave at the top and bottom right corner
    \coordinate (dc1) at ($(fig.north east) + (-0.8mm,-0.4mm)$);
    \coordinate (dc2) at ($(fig.south east) + (-0.8mm,0.4mm)$);
    \fill (dc1) circle (0.6pt) (dc2) circle (0.6pt);
    \draw[volt] ($(dc1)+(0,-1mm)$) -- ($(dc2)+(0,1mm)$);
    \node[anchor=west, color=gwvoltage] at ($(dc1)!0.5!(dc2) + (1.5mm,0)$)
      {$U_{dc}$};
  \end{tikzpicture}
\end{frame}

% ---- Converter operating modes -------------------------------------
\begin{frame}{Converter operating modes}
  \noindent The same converter works as rectifier or inverter, set by the direction of $P$.
  \vspace{1em}

  \centering
  \begin{tikzpicture}[node distance=0pt]
    \vscmode{C}{0mm,0mm}{Converter}{pboth}
    \vscmode{I}{-34mm,-30mm}{Inverter}{pdc2ac}
    \vscmode{R}{34mm,-30mm}{Rectifier}{pac2dc}

    \node[font=\scriptsize, color=atGray, anchor=east] at ($(C.west) + (-1mm,0)$) {AC};
    \node[font=\scriptsize, color=atGray, anchor=west] at ($(C.east) + (1mm,0)$) {DC};
  \end{tikzpicture}
\end{frame}


\begin{frame}{DC load flow definitions}
  \centering
  \includegraphics[width=0.8\textwidth]{\figfolder smp_level_1_label}
  \vspace{1em}
  \begin{itemize}
    \item DC equipment (incl. DC line) voltage is defined by $U_{dpe\_rec}$
    \item Converter size is defined by $U_{dpp\_rec}$
  \end{itemize}
\end{frame}


\begin{frame}{Monopole}
  \begin{columns}[c,onlytextwidth]
    \begin{column}{0.6\textwidth}\centering
      \topoimg[width=\linewidth,height=0.72\textheight,keepaspectratio]{\figfolder monopole}
    \end{column}
    \begin{column}{0.40\textwidth}
      \textbf{DC load flow}
      \begin{align*}
        U_{dpp\_rec} &= U_{dpe\_rec}\\
        I_{DC} &= \tfrac{U_{dpp\_rec}-U_{dpp\_inv}}{2 \cdot R_{DC}}\\
        U_{dne\_rec} &= I_{d}\cdot R_{DC}
      \end{align*}
      \begin{itemize}
        \item Max.\ voltage created by converter: $U_{dpp\_rec}$
        \item Max. equipment voltage: 
        \item \textcolor{red}{DC cables have different voltage ratings}
      \end{itemize}
    \end{column}
  \end{columns}
\end{frame}

\begin{frame}{Monopole - extending power transmission}
  \begin{columns}[c,onlytextwidth]
    \begin{column}{0.4\textwidth}
      \begin{itemize}
        \setlength{\itemsep}{0.3em}
        \item \textbf{Higher DC voltage}
            \begin{itemize}
              \item Limited equipment availability
              \item No redundancy
            \end{itemize}
        \item \textbf{Higher DC current}
            \begin{itemize}
              \item Higher DC line costs
              \item No redundancy
            \end{itemize}
        \item \textbf{Additional monopole}
            \begin{itemize}
              \item Additional converter equipment
              \item 4 DC conductors
              \item Increased reliability
            \end{itemize}
      \end{itemize}
    \end{column}
    \begin{column}{0.6\textwidth}\centering
      \topoimg[width=\linewidth,height=0.34\textheight,keepaspectratio]{\figfolder monopole}\\[0.4em]
      \topoimg[width=\linewidth,height=0.34\textheight,keepaspectratio]{\figfolder monopole}
    \end{column}
  \end{columns}
\end{frame}

\begin{frame}{Monopole $\rightarrow$ Bipole (BIP)}
  \centering
  \begin{columns}[c,onlytextwidth]
    \begin{column}{0.42\textwidth}\centering
      \topoimg[width=\linewidth,height=0.5\textheight,keepaspectratio]{\figfolder monopole_neg}\\[0.4em]
      \topoimg[width=\linewidth,height=0.5\textheight,keepaspectratio]{\figfolder monopole_neg}
    \end{column}
    \begin{column}{0.06\textwidth}\centering $\Rightarrow$ \end{column}
    \begin{column}{0.42\textwidth}\centering
      \topoimg[width=\linewidth,height=0.8\textheight,keepaspectratio]{\figfolder bip}
    \end{column}
  \end{columns}
\end{frame}

\begin{frame}{Bipole (BIP)}
  \begin{columns}[c,onlytextwidth]
    \begin{column}{0.54\textwidth}
      \begin{itemize}
        \setlength{\itemsep}{0.4em}
        \item 3 conductors of 2 different types.
        \item Two poles that can transmit power independently.
        \item On a fault in one pole, the other pole still transmits power ---
              redundancy.
        \item In symmetrical operation (both poles equal), no current flows
              through the DMR.
      \end{itemize}
    \end{column}
    \begin{column}{0.42\textwidth}\centering
      \topoimg[width=\linewidth,height=0.66\textheight,keepaspectratio]{\figfolder bip_conductor}
    \end{column}
  \end{columns}
\end{frame}

\begin{frame}{BIP operation modes}
  \begin{columns}[T,onlytextwidth]
    \begin{column}{0.32\textwidth}\centering
      \topoimg[width=\linewidth,height=0.42\textheight,keepaspectratio]{\figfolder bip_balanced}\\[0.3em]
      {\footnotesize\textbf{Bipolar --- balanced}\\ 
      \begin{itemize}
        \item Equal power for poles
        \item No DMR current
        \item Loss-optimized
      \end{itemize}
      }
    \end{column}
    \begin{column}{0.32\textwidth}\centering
      \topoimg[width=\linewidth,height=0.42\textheight,keepaspectratio]{\figfolder bip_monopolar}\\[0.3em]
      {\footnotesize\textbf{Monopolar}\\
      \begin{itemize}
        \item Emergency mode
        \item System-design definition
      \end{itemize}
      }
    \end{column}
    \begin{column}{0.32\textwidth}\centering
      \topoimg[width=\linewidth,height=0.42\textheight,keepaspectratio]{\figfolder bip_unbalanced}\\[0.3em]
      {\footnotesize\textbf{Bipolar --- unbalanced}\\
      \begin{itemize}
        \item Split-busbar operation
        \item DMR current flow
      \end{itemize}
      }
    \end{column}
  \end{columns}
\end{frame}

% ---- Symmetrical Monopole -----------------------------------
\begin{frame}{Symmetrical Monopole (SMP)}
  \begin{columns}[c,onlytextwidth]
    \begin{column}{0.54\textwidth}\centering
      \topoimg[width=\linewidth,height=0.72\textheight,keepaspectratio]{\figfolder smp}
    \end{column}
    \begin{column}{0.42\textwidth}
      \begin{itemize}
        \setlength{\itemsep}{0.25em}
        \item High-impedance earth reference on the AC side
        \item $U_{dpp\_Rec}= 2 \cdot \lvert U_{dpe\_Rec} \rvert$
        \item $I_d=\frac{U_{dpp\_rec} - U_{dpp\_inv}}{2 \cdot R_d}$
        \item Only for VSC-HVDC systems
        \item \textcolor{green}{DC cables have the same voltage rating} 
      \end{itemize}
    \end{column}
  \end{columns}
\end{frame}

% ---- MMC comparison table -----------------------------------
\begin{frame}{Modular Multilevel Converter (MMC)}
  Power transmission parameters:
  \begin{itemize}
    \item Power: 1000\,MW
    \item DC voltage: $\pm$ 320 kV
    \item Max. DC current: 1563 A
  \end{itemize}
  \vspace{0.6em}

  \centering
  \renewcommand{\arraystretch}{1.3}
  \begin{tabular}{@{}l c c c@{}}
     & \textbf{2 $\times$ Monopole} & \textbf{1 $\times$ BIP} & \textbf{1 $\times$ SMP} \\
    \hline
    Transformers      & \textcolor{red}{$--$} & \textcolor{red}{$--$} & \textcolor{green}{$++$} \\
    DC lines          & \textcolor{red}{4 lines / 2 types} & \textcolor{red}{3 lines / 2 types} & \textcolor{green}{2 lines / 1 type}\\
    Converters        & \textcolor{red}{4 converters} & \textcolor{red}{4 converters} & \textcolor{green}{2 converters}\\
    Converter rating ($U_{dpp}$) & \textcolor{green}{320\,kV} & \textcolor{green}{320\,kV} & \textcolor{red}{640\,kV} \\
    Losses            & \textcolor{red}{$--$} & \textcolor{red}{$-$} & \textcolor{green}{$+$} \\
    Reliability       & \textcolor{red}{$++$} & \textcolor{red}{$+$} & \textcolor{red}{$--$} \\
    \hline
  \end{tabular}
\end{frame}

% ---- DC cable -----------------------------------------------
\begin{frame}{DC cable}
  \begin{columns}[c,onlytextwidth]
    \begin{column}{0.56\textwidth}
      \begin{itemize}
        \setlength{\itemsep}{0.25em}
        \item \textbf{Insulator} --- XLPE (cross-linked polyethylene) or
              MI (mass-impregnated).
        \item \textbf{Conductor} --- copper or aluminium.
        \item Balance between DC voltage and DC current for the most
              economical design.
        \item Insulator is generally cheaper than conductor
              $\rightarrow$ higher voltage rating is preferred
      \end{itemize}
      \vspace{0.3em}
      Project voltages: 80/200/250/\textbf{320}/400/\textbf{525}\,kV\\
      (525\,kV is new standard for $\ge 2000$\,MW).
    \end{column}
    \begin{column}{0.40\textwidth}\centering
      \topoimg[width=\linewidth,height=0.4\textheight,keepaspectratio]{\figfolder cable_crosssection}\\[0.4em]
      \topoimg[width=0.8\linewidth,height=0.28\textheight,keepaspectratio]{\figfolder cable}\\
      {\color{atGray}\scriptsize www.borealisgroup.com}
    \end{column}
  \end{columns}
\end{frame}


% ---- Topology comparison table ------------------------------
\begin{frame}{Topology comparison}
  \centering
  \renewcommand{\arraystretch}{1.3}
  \begin{tabular}{@{}p{4.6cm} c c c@{}}
     & \textbf{Sym.\ Monopole} & \textbf{Rigid Bipole} & \textbf{Bipole w/ MR} \\
    \hline
    DC conductors                         & 2 & 2 & 3 \\
    Power on converter fault              & 0\% & 50\% & 50\% \\
    Power on DC conductor fault           & 0\% & 0\% & 50\% \\
    Economical DC voltage                 & $\le 400$\,kV & $>500$\,kV & $>500$\,kV \\
    Costs                                 & \$ & \$\$ & \$\$\$ \\
    Max.\ power                           & \makecell{1400\,MW (320\,kV)\\1500\,MW (400\,kV)} & 2000\,MW & 3000\,MW \\
    \hline
  \end{tabular}
\end{frame}


\end{document}
